% ══════════════════════════════════════════════════════════════════
% Section 1 — Introduction
% ══════════════════════════════════════════════════════════════════
\section{Introduction}
\label{sec:introduction}

Nonlinear finite element simulations of soft materials require the repeated
evaluation of constitutive relations --- strain energy functions (SEFs) and
their derivatives --- at every integration point and every Newton iteration.
For complex hyperelastic models, especially those involving fiber-reinforced
anisotropy, the evaluation of stress tensors and consistent tangent operators
constitutes a significant fraction of the total computational cost.  This
bottleneck motivates the development of \emph{data-driven surrogates}: neural
networks trained to approximate the constitutive response and deployable
inside existing finite element solvers.

\subsection{Related work}

Early neural network surrogates for constitutive modeling employed standard
feedforward networks trained on stress--strain data
\cite{% TODO: Ghaboussi1991, Hashash2004, Lefik2003
}.
While effective for interpolation, unconstrained networks can violate
fundamental thermodynamic requirements such as polyconvexity, objectivity, and
the stress-free reference configuration, leading to non-physical predictions
and convergence failures in implicit solvers.

Physics-informed architectures address these limitations.  Input-convex neural
networks (ICNNs) \cite{% TODO: Amos2017
} guarantee convexity of the learned energy with respect to its inputs, a
necessary condition for material stability.  Klein et al.~\cite{%
TODO: Klein2022
} demonstrated polyconvex ICNNs for hyperelasticity by decomposing the energy
into convex functions of separate invariant groups.  Constitutive artificial
neural networks (CANNs) \cite{% TODO: Linka2023
} impose structure through predefined basis functions with learnable weights,
enabling model discovery and interpretability.

Despite these advances, several gaps remain:
\begin{itemize}
  \item \textbf{Systematic comparison.}  Most studies evaluate a single
    architecture on one or two materials.  A fair comparison of MLP, ICNN,
    and polyconvex ICNN across diverse isotropic and anisotropic materials
    under identical training conditions is lacking.
  \item \textbf{Energy-based training with stress accuracy.}  Training on
    energy alone does not guarantee accurate stress predictions; training on
    stress alone loses the energy potential.  The joint energy--stress loss via
    automatic differentiation is underexplored for convexity-preserving
    architectures.
  \item \textbf{Deployment in production solvers.}  Trained models typically
    remain in Python, disconnected from Fortran-based FE codes.  Automated
    export of neural networks into standalone UMAT subroutines --- including
    analytical stress and tangent computation --- has not been systematically
    addressed.
  \item \textbf{Finite element validation.}  End-to-end validation of
    NN-based UMATs against analytical solutions in commercial solvers
    (Abaqus) and research codes (FEAP) is rare.
\end{itemize}

\subsection{Contributions}

This paper makes the following contributions:

\begin{enumerate}
  \item A systematic benchmark of three neural architectures (MLP, ICNN,
    polyconvex ICNN) on five hyperelastic materials (three isotropic, two
    anisotropic), using an energy--stress dual loss with automatic
    differentiation (\Cref{sec:results}).
  \item Scaling studies quantifying the effect of training set size and
    network width on surrogate accuracy (\Cref{sec:results:scaling}).
  \item An automated Fortran~90 export pipeline that emits self-contained
    UMAT subroutines with embedded network weights, analytical kinematics,
    and consistent tangent operators (\Cref{sec:fortran_export}).
  \item Finite element validation of the exported hybrid NN--UMAT against
    analytical UMATs in Abaqus and FEAP under uniaxial, biaxial, and shear
    loading (\Cref{sec:fe_validation}).
  \item An open-source Python library, \texttt{hyper-surrogate}, that
    implements the entire pipeline from data generation to Fortran code
    emission.
\end{enumerate}

\subsection{Outline}

\Cref{sec:continuum} reviews the continuum mechanics framework and the
material models considered.  \Cref{sec:architectures} describes the neural
network architectures.  \Cref{sec:training} details the training procedure and
the energy--stress loss.  \Cref{sec:fortran_export} presents the Fortran
export pipeline.  \Cref{sec:results} reports benchmark results, and
\Cref{sec:fe_validation} presents the finite element validation.
\Cref{sec:discussion} discusses limitations and future directions, and
\Cref{sec:conclusion} concludes the paper.
